% ─── CAPA ───
\begin{capa}
    \centering
    \vspace*{4cm}
    {\ABNTEXchapterfont\Large\bfseries OpenCode Ecosystem v5.4.0 R23}\\[1cm]
    {\ABNTEXchapterfont\large Arquitetura, Implementação e Validação de um Sistema de\\Engenharia de Software com Metacognição Funcional\\e Behavioral Gate Preventivo}\\[3cm]
    {\large Marcelo Claro Laranjeira}\\[4cm]
    {\large Crateús, Ceará\\2026}
\end{capa}

% ─── FOLHA DE ROSTO ───
\begin{folhaderosto}
    \centering
    {\ABNTEXchapterfont\Large Marcelo Claro Laranjeira}\\[1.5cm]
    {\ABNTEXchapterfont\Large\bfseries OpenCode Ecosystem v5.4.0 R23}\\[2cm]
    \begin{flushright}
    \begin{minipage}{0.5\textwidth}
    Dissertação apresentada como requisito para documentação científica do projeto OpenCode Ecosystem. Gerada pelo próprio sistema (AutoEvolve v5.4.0), criado e orientado pelo pesquisador Marcelo Claro Laranjeira.
    \end{minipage}
    \end{flushright}
    \vfill
    {\large Crateús, Ceará\\2026}
\end{folhaderosto}

% ─── FOLHA DE APROVAÇÃO ───
\begin{folhadeaprovacao}
    \centering
    {\ABNTEXchapterfont\Large Marcelo Claro Laranjeira}\\[1cm]
    {\ABNTEXchapterfont\large\bfseries OpenCode Ecosystem v5.4.0 R23}\\[2cm]
    \begin{flushright}
    \begin{minipage}{0.5\textwidth}
    Dissertação validada por 312 testes críticos (15 suites TDD, 100\% de aprovação) e auditada pelo próprio sistema via SPEC-036 MetacognitiveMonitor e SPEC-038 TrustEngine.
    \end{minipage}
    \end{flushright}
    \vspace{1cm}
    Aprovado em: \_\_\_ / \_\_\_ / 2026
\end{folhadeaprovacao}

% ─── DEDICATÓRIA ───
\begin{dedicatoria}
    \vspace*{8cm}
    \begin{flushright}
    Àqueles que acreditam que sistemas de software podem\\
    não apenas executar tarefas, mas também observar,\\
    corrigir, aprender e prevenir — a si mesmos.
    \end{flushright}
\end{dedicatoria}

% ─── AGRADECIMENTOS ───
\begin{agradecimentos}
Agradeço ao ecossistema OpenCode por ter se auto-diagnosticado, identificado seus próprios gaps, e implementado as correções necessárias — incluindo a capacidade de gerar esta dissertação como demonstração de metacognição funcional. Agradeço aos 312 testes críticos que garantem que cada afirmação neste documento é verificável e reproduzível. Agradeço aos projetos open-source que inspiraram componentes específicos: Elinor Ostrom (DP1-DP8), George Pólya (estratégias de resolução), Atkinson \& Shiffrin (memória com esquecimento natural), e Karpathy (autoresearch loops).
\end{agradecimentos}

% ─── EPÍGRAFE ───
\begin{epigrafe}
    \vspace*{8cm}
    \begin{flushright}
    \textit{``Um sistema que identifica seus próprios gaps mas não pode corrigi-los\\é um sistema incompleto. Um sistema que corrige mas não previne\\é um sistema reativo. A metacognição completa exige ambos.''}\\
    \medskip
    --- MetacognitiveMonitor, SPEC-036, OpenCode Ecosystem (2026)
    \end{flushright}
\end{epigrafe}

% ─── RESUMO ───
\begin{resumo}
Esta dissertação apresenta o OpenCode Ecosystem v5.4.0 R23, uma plataforma de engenharia de software com agentes inteligentes que implementa um pipeline completo de scanners epistemológicos, decomposição de conhecimento em insumos cognitivos, metacognição funcional (N3.5), e behavioral gate preventivo com trust scoring adaptativo. O sistema é composto por 22 módulos Python (8.937 linhas), validados por 15 suites TDD totalizando 312 testes críticos com 100\% de aprovação e zero regressão. A arquitetura é organizada em 7 camadas (Skills $\rightarrow$ MCPs $\rightarrow$ Agentes $\rightarrow$ Scanner Pipeline $\rightarrow$ MCSP $\rightarrow$ Metacognição $\rightarrow$ Trust Engine), documentada por 10 ADRs e 13 especificações formais (SPEC-025 a SPEC-038). O sistema completou 23 ciclos evolutivos documentados, progredindo de um score inicial de 85/100 para 100/100. A contribuição central é a demonstração de que metacognição simbólica funcional — auto-observação, detecção de anomalias com thresholds adaptativos, correção automática com aprendizado de eficácia, inferência causal de causas raiz, e prevenção baseada em confiança — é implementável em um sistema de engenharia de software real, com todas as decisões rastreáveis e todos os testes reproduzíveis.

\textbf{Palavras-chave:} OpenCode Ecosystem. Metacognição Artificial. Behavioral Gate. Trust Scoring. Test-Driven Development. Engenharia de Software com Agentes. Composição Unitária do Conhecimento. Governança Cooperativa (Ostrom). N3.5.
\end{resumo}

% ─── ABSTRACT ───
\begin{resumo}[Abstract]
This dissertation presents the OpenCode Ecosystem v5.4.0 R23, a software engineering platform with intelligent agents implementing a complete pipeline of epistemological scanners, knowledge decomposition into cognitive inputs, functional metacognition (N3.5), and a preventive behavioral gate with adaptive trust scoring. The system comprises 22 Python modules (8,937 lines), validated by 15 TDD suites totaling 312 critical tests with 100\% pass rate and zero regression. The central contribution is the demonstration that functional symbolic metacognition — self-observation, anomaly detection with adaptive thresholds, automatic correction with efficacy learning, causal root cause inference, and trust-based prevention — is implementable in a real software engineering system, with all decisions traceable and all tests reproducible.

\textbf{Keywords:} OpenCode Ecosystem. Artificial Metacognition. Behavioral Gate. Trust Scoring. TDD. Agent-Based Software Engineering. Unit Composition of Knowledge. Cooperative Governance (Ostrom). N3.5.
\end{resumo}

% ─── LISTAS ───
\pdfbookmark[0]{\listfigurename}{lof}
\listoffigures*
\cleardoublepage
\pdfbookmark[0]{\listtablename}{lot}
\listoftables*
\cleardoublepage
\pdfbookmark[0]{\contentsname}{toc}
\tableofcontents*
\cleardoublepage
